\chapter{Gruppenaxiome und unmittelbare Folgerungen}

Dieser Leitfaden verzeichnet zentrale Definitionen, Charakterisierungen und
Hauptsätze für beliebige Gruppen. Die Ausgangsgruppe wird weder als endlich
noch als abelsch, endlich erzeugt, auflösbar, geordnet oder topologisch
vorausgesetzt. Bedingungen an einzelne Elemente, Untergruppen oder
Homomorphismen stehen ausdrücklich in der jeweiligen Aussage.

Aufgenommen sind Ergebnisse, für deren Aussage die Gruppenstruktur
wesentlich ist. Allgemeine Sätze über Funktionen, Injektivität,
Surjektivität, Bijektivität, Äquivalenzrelationen, Halbordnungen,
totale Ordnungen, Schranken, Infima und Suprema werden nicht wiederholt.
Die entsprechenden Begriffe werden bei Bedarf verwendet. Auch die
allgemeine Halbgruppen- und Monoidtheorie sowie die Spezialtheorien
abelscher und endlicher Gruppen bleiben ihren eigenen Bänden vorbehalten.
Gruppenspezifische Konstruktionen dürfen neue Gruppen mit besonderen
Eigenschaften liefern; solche Eigenschaften werden der Ausgangsgruppe
nicht zusätzlich auferlegt.

Die Auswahl enthält keine Beweise oder Beweisskizzen und erhebt keinen
Anspruch auf ein vollständiges Verzeichnis aller bekannten Gruppensätze.
Sie erweitert den bisher in Band~40 ausgearbeiteten Bestand um klassische
Literaturresultate. Als mengentheoretische Grundlage gilt ZFC einschließlich
des Auswahlaxioms. Die Satznummern mit Präfix \(40H\) gehören ausschließlich
zu diesem Leitfaden.

\section{Grundbegriff und Notation}

\begin{definition}[Gruppe]
Eine \emph{Gruppe} ist eine Menge \(G\) mit einer inneren Verknüpfung
\((a,b)\mapsto ab\), einem Element \(1\in G\) und folgenden Eigenschaften:
\[
 (ab)c=a(bc),\qquad 1a=a1=a
 \quad(a,b,c\in G),
\]
und zu jedem \(a\in G\) gibt es ein \(b\in G\) mit \(ab=ba=1\).
Insbesondere ist \(G\) nichtleer. Im ganzen Leitfaden bezeichnen
\(G,H,K\) Gruppen, soweit nichts anderes angegeben wird.
Die Gruppe mit genau einem Element heißt \emph{trivial}.
\end{definition}

\begin{definition}[Gruppenhomomorphismus und Isomorphismus]
Eine Abbildung \(f:G\to H\) heißt \emph{Gruppenhomomorphismus}, wenn
\(f(ab)=f(a)f(b)\) für alle \(a,b\in G\) gilt.
Ein bijektiver Gruppenhomomorphismus heißt \emph{Isomorphismus};
man schreibt \(G\cong H\), wenn ein solcher existiert.
Ein injektiver Homomorphismus heißt \emph{Einbettung}.
Es seien
\[
 \ker f=\{g\in G:f(g)=1_H\},\qquad
 \operatorname{im}f=f(G).
\]
Ein Homomorphismus \(G\to G\) heißt \emph{Endomorphismus},
ein Isomorphismus \(G\to G\) \emph{Automorphismus}.
\end{definition}

\begin{theorem}[Eindeutigkeit von Einheit und Inversen]
Das neutrale Element einer Gruppe ist eindeutig. Zu jedem \(a\in G\)
gibt es genau ein beidseitiges inverses Element, bezeichnet mit \(a^{-1}\).
Schon \(ab=1\) oder \(ba=1\) erzwingt \(b=a^{-1}\).
Es gelten
\[
 1^{-1}=1,\qquad (a^{-1})^{-1}=a,\qquad
 (ab)^{-1}=b^{-1}a^{-1}.
\]
Für \(n\geq1\) ist
\((a_1\cdots a_n)^{-1}=a_n^{-1}\cdots a_1^{-1}\).
\cite{project40,milne}
\end{theorem}

\begin{theorem}[Kürzung und eindeutige Lösung von Gleichungen]
Für \(a,b,c\in G\) gelten
\[
 ab=ac\iff b=c,\qquad ba=ca\iff b=c.
\]
Die Gleichungen \(ax=b\), \(ya=b\) und \(axb=c\) haben jeweils genau eine
Lösung in \(G\), nämlich
\[
 x=a^{-1}b,\qquad y=ba^{-1},\qquad x=a^{-1}cb^{-1}.
\]
\cite{project40,milne}
\end{theorem}

\begin{theorem}[Keine weiteren idempotenten Elemente]
Für \(a\in G\) gilt \(a^2=a\) genau dann, wenn \(a=1\).
Ein Element \(z\in G\) mit \(zg=gz=z\) für alle \(g\in G\)
existiert genau dann, wenn \(G\) trivial ist. \cite{project40}
\end{theorem}

\section{Charakterisierungen der Gruppenstruktur}

\begin{theorem}[Einseitige Axiome genügen]
Es sei \(G\ne\varnothing\) mit einer assoziativen inneren Verknüpfung.
Gibt es ein \(e\in G\) mit \(ea=a\) für alle \(a\in G\), und besitzt jedes
\(a\in G\) ein \(b\in G\) mit \(ba=e\), so ist \(G\) eine Gruppe mit
Einheit \(e\). Ebenso genügt die gemeinsame rechtsseitige Variante
\(ae=a\) und \(ab=e\).
Die jeweils einseitigen Voraussetzungen müssen auf derselben Seite stehen.
\cite{une}
\end{theorem}

\begin{theorem}[Charakterisierung durch Division]
\label{b40guide:thm:division}
Für eine nichtleere Menge \(G\) mit assoziativer innerer Verknüpfung
sind äquivalent:
\begin{enumerate}
\item \(G\) ist eine Gruppe.
\item Für alle \(a,b\in G\) sind \(ax=b\) und \(ya=b\) in \(G\) lösbar.
\item Für alle \(a,b\in G\) sind diese beiden Gleichungen eindeutig lösbar.
\end{enumerate}
In Bedingung~2 muß die Eindeutigkeit also nicht zusätzlich verlangt werden.
\cite{une}
\end{theorem}

\begin{theorem}[Translationscharakterisierung]
Auf einer nichtleeren Menge mit assoziativer innerer Verknüpfung
definiere man \(L_a(x)=ax\) und \(R_a(x)=xa\).
Die Verknüpfung ist genau dann eine Gruppenverknüpfung, wenn alle
\(L_a\) und \(R_a\) Permutationen der Menge sind.
Für eine Gruppe gelten
\[
 L_a^{-1}=L_{a^{-1}},\quad R_a^{-1}=R_{a^{-1}},
 \quad L_aR_b=R_bL_a.
\]
Hier bezeichnet die Verkettung \(uv\) stets \(u\circ v\).
\cite{milne}
\end{theorem}

\section{Ganzzahlige Potenzen und Elementordnungen}

\begin{definition}[Potenzen und Ordnung eines Elements]
Es seien \(a^0=1\), \(a^{n+1}=a^na\) für \(n\geq0\) und
\(a^{-n}=(a^{-1})^n\). Die \emph{Ordnung} \(\operatorname{ord}(a)\)
ist die kleinste positive Zahl \(d\) mit \(a^d=1\), falls eine solche
existiert, und sonst \(\infty\). Diese Ordnung ist eine Zahl bzw. das
Symbol \(\infty\), keine Ordnungsrelation auf \(G\).
\end{definition}

\begin{theorem}[Ganzzahlige Potenzgesetze]
Für \(a\in G\) und \(m,n\in\mathbb Z\) gelten
\[
 a^{m+n}=a^ma^n,\qquad (a^m)^n=a^{mn},\qquad
 (a^m)^{-1}=a^{-m}.
\]
Für kommutierende Elemente \(a,b\in G\), also \(ab=ba\), gilt zusätzlich
\((ab)^n=a^nb^n\) für jedes \(n\in\mathbb Z\).
Eine Kommutativität der gesamten Gruppe ist hierzu nicht erforderlich.
\cite{milne}
\end{theorem}

\begin{theorem}[Vollständige Alternative für die Potenzen eines Elements]
Für jedes \(a\in G\) tritt genau einer der folgenden Fälle ein:
\begin{enumerate}
\item \(\operatorname{ord}(a)=\infty\), und für \(m,n\in\mathbb Z\) gilt
\(a^m=a^n\iff m=n\).
\item \(\operatorname{ord}(a)=d<\infty\), und für \(m,n\in\mathbb Z\) gilt
\(a^m=a^n\iff d\mid(m-n)\).
\end{enumerate}
Im zweiten Fall ist für jedes \(k\in\mathbb Z\)
\[
 \operatorname{ord}(a^k)=\frac{d}{\gcd(d,|k|)}
 \quad\text{mit }\gcd(d,0)=d.
\]
Im ersten Fall hat \(a^k\) für \(k\ne0\) unendliche Ordnung.
\cite{milne}
\end{theorem}

\begin{theorem}[Konjugation erhält Potenzen und Elementordnungen]
Für \(a,g\in G\) und \(n\in\mathbb Z\) gelten
\[
 (gag^{-1})^n=ga^ng^{-1},\qquad
 \operatorname{ord}(gag^{-1})=\operatorname{ord}(a).
\]
Auch \(\operatorname{ord}(a^{-1})=\operatorname{ord}(a)\).
\cite{milne}
\end{theorem}
